\chapter{Conception et architecture}
\label{chap:conception}

\section*{Introduction}
\addcontentsline{toc}{section}{Introduction}

Ce chapitre présente l'architecture retenue et justifie les choix technologiques. L'accent est mis
sur les décisions structurantes --- celles qui, une fois prises, contraignent tout le reste --- et
sur les arbitrages assumés, y compris les trois qui restent ouverts et sont énoncés comme tels en
fin de chapitre.

\section{Architecture générale}

Le système se décompose en deux pipelines distincts : un pipeline d'\textbf{indexation} exécuté
une seule fois hors ligne, et un pipeline de \textbf{requête} exécuté à chaque question
(figure~\ref{fig:architecture}).

\begin{figure}[htbp]
  \centering
  % Architecture générale.
%
% Principes de tracé, tenus strictement :
%   - flux principal vertical, une seule colonne : la requête se lit de haut
%     en bas sans jamais revenir en arrière ;
%   - la chaîne hors ligne occupe une colonne parallèle et ne rejoint la
%     colonne principale qu'en un seul point — l'étape de recherche, qui est
%     précisément le seul endroit où les deux pipelines se rencontrent ;
%   - routage orthogonal exclusivement, aucun segment oblique ;
%   - aucun croisement d'arêtes ;
%   - gabarits de nœuds uniformes par colonne.
\begin{tikzpicture}[
    font=\sffamily\scriptsize,
    boite/.style={draw, line width=0.7pt, rounded corners=2pt,
                  text width=34mm, minimum height=9.5mm, align=center,
                  inner sep=3pt},
    principal/.style={boite, minimum height=10.5mm},
    horsligne/.style={boite, text width=30mm, minimum height=8.5mm},
    d/.style={draw=slate,     fill=mist,          text=ink},          % donnée
    t/.style={draw=inptblue,  fill=inptblue!8,    text=navy},         % traitement
    m/.style={draw=leafgreen, fill=leafgreen!8,   text=leafgreen!40!black}, % modèle
    stop/.style={draw=anrtred, fill=anrtred!6,    text=anrtred!70!black},
    choix/.style={draw=amberdark, fill=amberdark!10, text=amberdark!55!black,
                  line width=0.7pt, diamond, aspect=2.3, align=center,
                  inner sep=1pt, text width=19mm},
    fleche/.style={-{Stealth[length=2.3mm,width=1.7mm]}, line width=0.7pt,
                   draw=slate, rounded corners=2.5pt},
    remise/.style={fleche, draw=inptblue!75, line width=0.9pt},
    refus/.style={fleche, draw=anrtred},
    tag/.style={font=\sffamily\scriptsize, text=slate, inner sep=2pt},
  ]

% ============================================================ grille
\def\XA{0}      % colonne principale (requête)
\def\XB{5.6}    % colonne hors ligne (indexation)
\def\XR{-4.5}   % dérivation d'abstention

% ============================================ colonne principale : la requête
\node[principal, d] (q)      at (\XA, 0)      {Question du citoyen};
\node[principal, m] (vec)    at (\XA,-1.75)   {Vectorisation\\BGE-M3};
\node[principal, t] (rdense) at (\XA,-3.5)    {Recherche dense};
\node[principal, t] (rlex)   at (\XA,-4.75)   {Recherche lexicale};
\node[principal, t] (rrf)    at (\XA,-6.4)    {Fusion RRF};
\node[choix]        (gate)   at (\XA,-8.05)   {score $\geq$\\\SeuilAbstention{} ?};
\node[principal, m] (gen)    at (\XA,-9.9)    {Génération\\Mistral 7B via Ollama};
\node[principal, t] (verif)  at (\XA,-11.65)  {Contrôle des citations};
\node[principal, d] (rep)    at (\XA,-13.4)   {Réponse citée\\+ avertissement};

\draw[fleche] (q)      -- (vec);
\draw[fleche] (vec)    -- (rdense);
\draw[fleche] (rlex)   -- (rrf);
\draw[fleche] (rrf)    -- (gate);
\draw[fleche] (gate)   -- node[tag, right=1pt, pos=0.35] {oui} (gen);
\draw[fleche] (gen)    -- (verif);
\draw[fleche] (verif)  -- (rep);

% Les deux recherches sont menées en parallèle sur la même question : la
% dérivation part du même point et les rejoint toutes deux.
\draw[fleche] (vec.west) -- ++(-1.05,0) |- (rlex.west);

% =================================== colonne hors ligne : la chaîne d'indexation
% Les deux index sont posés exactement sur les ordonnées des deux recherches
% qu'ils alimentent : la remise se fait par un segment horizontal pur.
\node[horsligne, d] (pdf)   at (\XB, 1.9)   {PDF officiel\\Loi 65-99};
\node[horsligne, t] (pars)  at (\XB, 0.5)   {\module{parser.py}\\extraction, découpage};
\node[horsligne, d] (jsonl) at (\XB,-0.9)   {Corpus JSONL\\\NbArticles{} articles};
\node[horsligne, m] (emb)   at (\XB,-2.3)   {Vectorisation\\BGE-M3};
\node[horsligne, d] (ichr)  at (\XB,-3.5)   {Index dense\\Chroma};
\node[horsligne, d] (ibm)   at (\XB,-4.75)  {Index lexical\\BM25};

\draw[fleche] (pdf)   -- (pars);
\draw[fleche] (pars)  -- (jsonl);
\draw[fleche] (jsonl) -- (emb);
\draw[fleche] (emb)   -- (ichr);
% BM25 indexe le texte, pas les vecteurs : il dérive du corpus, pas du modèle.
\draw[fleche] (jsonl.east) -- ++(1.15,0) |- (ibm.east);

% ============================ point de rencontre unique entre les deux chaînes
\draw[remise] (ichr.west) -- (rdense.east);
\draw[remise] (ibm.west)  -- (rlex.east);

% ================================================ dérivation d'abstention
\node[principal, stop] (abst) at (\XR,-8.05) {Abstention immédiate\\le modèle n'est pas appelé};
\draw[refus] (gate.west) -- node[tag, above, text=anrtred] {non} (abst.east);

% Aucune annotation en prose dans la figure : le commentaire appartient à la
% légende et au corps du texte, pas au schéma.

% Cadres de regroupement, posés en arrière-plan.
\begin{scope}[on background layer]
  \node[draw=rule, line width=0.7pt, rounded corners=3pt, dash pattern=on 3pt off 2pt,
        inner xsep=4mm, inner ysep=5mm, fit=(pdf)(pars)(jsonl)(emb)(ichr)(ibm)] (bhl) {};
  % L'abstention fait partie du pipeline de requête : le cadre l'englobe.
  \node[draw=rule, line width=0.7pt, rounded corners=3pt, dash pattern=on 3pt off 2pt,
        inner xsep=4mm, inner ysep=5mm,
        fit=(q)(vec)(rdense)(rlex)(rrf)(gate)(gen)(verif)(rep)(abst)] (ben) {};
\end{scope}

\node[font=\sffamily\scriptsize\bfseries, text=slate, anchor=south, align=center]
  at (bhl.north) {INDEXATION\\[-1pt]\scriptsize\mdseries hors ligne, une seule fois};
% Le cadre de requête étant élargi par la dérivation d'abstention, son titre
% se cale sur la colonne principale plutôt que sur le milieu du cadre.
\node[font=\sffamily\scriptsize\bfseries, text=slate, anchor=south, align=center]
  at (q.north |- ben.north) {REQUÊTE\\[-1pt]\scriptsize\mdseries en ligne, à chaque question};

\end{tikzpicture}

  \caption[Architecture générale du système]{Architecture générale. La requête se lit de haut en
  bas ; la chaîne d'indexation, exécutée une seule fois hors ligne, occupe la colonne de droite.
  Les deux pipelines ne se rencontrent qu'en \emph{un seul point} --- la remise des deux index aux
  deux recherches (flèches bleues). Le modèle d'embedding apparaît des deux côtés : c'est
  volontairement le même, et c'est cette identité qui donne un sens à la comparaison entre une
  question et un article.}
  \label{fig:architecture}
\end{figure}

\medskip
\centerline{\legendeschema}
\medskip

Le modèle BGE-M3 intervient \textbf{deux fois} : à l'indexation, pour vectoriser les articles ; à
la requête, pour vectoriser la question. Il ne s'agit pas d'une commodité d'implémentation mais
d'une nécessité : deux modèles différents produiraient des vecteurs situés dans des espaces sans
rapport, et la similarité calculée entre eux n'aurait aucune signification.

\subsection{Deux principes structurants}

\paragraph{Le garde-fou d'abstention est placé \emph{avant} le modèle}
Il s'agit d'un test de seuil sur le score de recherche, donc d'un mécanisme
\textbf{indépendant du modèle}. Il ne dépend pas de la bonne volonté du LLM à dire «~je ne sais
pas~» --- ce sur quoi les petits modèles locaux se révèlent peu fiables, comme le
chapitre~\ref{chap:evaluation} le montrera concrètement. Ce choix a un effet secondaire mesurable
et non recherché : lorsqu'il se déclenche, le modèle n'est pas appelé du tout, et la réponse est
rendue en une fraction du temps habituel.

\paragraph{La recherche et la génération restent séparables}
À tout moment il est possible d'afficher les articles retrouvés \emph{avant} l'appel au modèle.
Face à une mauvaise réponse, cette séparation permet de déterminer immédiatement si le problème
vient de la recherche ou de la génération --- distinction impossible à faire si l'on n'observe que
la sortie finale. Elle rend aussi l'évaluation comparative possible : on peut mesurer la recherche
seule, sans que le modèle n'intervienne dans le résultat.

\section{Séquence d'une requête}

La figure~\ref{fig:sequence} détaille le trajet d'une question, en indiquant à chaque étape le
temps médian mesuré au chapitre~\ref{chap:evaluation}.

\begin{figure}[htbp]
  \centering
  % Séquence d'une requête, avec le temps médian mesuré à chaque étape.
% Les deux chemins possibles — réponse ou abstention — sont représentés
% ensemble : c'est leur écart de coût qui fait l'intérêt de la figure.
\begin{tikzpicture}[font=\sffamily\scriptsize, node distance=5.5mm]

\def\colx{0}          % colonne principale
\def\tempsx{4.55}     % gouttière des temps

% ------------------------------------------------------------------- étapes
\node[donnee, text width=34mm] (q) at (\colx,0) {Question reçue};
\node[modele, text width=34mm, below=of q] (vec)
  {Vectorisation de la question (BGE-M3)};
\node[traitement, text width=34mm, below=of vec] (rech)
  {Interrogation des deux index, fusion RRF};
\node[decision, text width=22mm, below=7mm of rech] (gate)
  {meilleur score $\geq$ \SeuilAbstention{} ?};
\node[modele, text width=34mm, below=9mm of gate] (gen)
  {Génération sous prompt strict, température 0,1};
\node[traitement, text width=34mm, below=of gen] (verif)
  {Extraction des numéros cités et confrontation au contexte fourni};
\node[donnee, text width=34mm, below=of verif] (out)
  {Réponse, articles sources, avertissement};

\draw[flux] (q) -- (vec);
\draw[flux] (vec) -- (rech);
\draw[flux] (rech) -- (gate);
\draw[flux] (gate) -- node[etiq, right, pos=0.35]{oui} (gen);
\draw[flux] (gen) -- (verif);
\draw[flux] (verif) -- (out);

% ---------------------------------------------------- branche d'abstention
\node[arret, text width=28mm, left=13mm of gen] (abst)
  {Phrase d'abstention rendue telle quelle};
\draw[rejet] (gate) -| node[etiq, above, pos=0.25, text=anrtred]{non} (abst);

\node[annot, below=4mm of abst, text width=32mm]
  {Le garde-fou est un test de seuil : il ne dépend pas de la bonne volonté
   du modèle à dire «\,je ne sais pas\,».};

% ------------------------------------------------------ gouttière des temps
\node[etiq, anchor=west, text=slate] at (\tempsx,{0.32}) {\bfseries médiane};
\node[etiq, anchor=west] at (\tempsx,{-1.55}) {\LatRecherche{} s};
\node[etiq, anchor=west, text width=27mm] at (\tempsx,{-2.2})
  {recherche complète};
\node[etiq, anchor=west] at (\tempsx,{-5.3}) {\LatGeneration{} s};
\node[etiq, anchor=west, text width=27mm] at (\tempsx,{-5.95}) {génération};
\node[etiqfort, anchor=west] at (\tempsx,{-8.6}) {\LatTotale{} s};
\node[etiq, anchor=west, text width=27mm] at (\tempsx,{-9.25}) {total};

\draw[rule, line width=0.5pt] ($(\tempsx-0.28,0.55)$) -- ($(\tempsx-0.28,-9.8)$);

% Le chemin d'abstention, chiffré. L'étiquette se place à gauche de la
% descente rouge, pour ne pas la chevaucher.
\node[etiq, text=anrtred, anchor=east, align=right]
  at ($(abst.north)+(-2.5mm,8mm)$)
  {\LatAbstention{} s\\\scriptsize soit \FacteurGarde{}$\times$ plus rapide};

\end{tikzpicture}

  \caption[Séquence d'une requête]{Séquence d'une requête, avec les deux issues possibles. La
  branche d'abstention court-circuite l'étape la plus coûteuse : c'est pourquoi une question hors
  périmètre reçoit sa réponse \FacteurGarde{} fois plus vite qu'une question traitée.}
  \label{fig:sequence}
\end{figure}

\section{Rôle de chaque module}

Le système compte cinq modules, chacun responsable d'une seule étape de la chaîne
(tableau~\ref{tab:modules}).

\begin{table}[htbp]
  \centering
  \caption{Répartition des responsabilités entre modules}
  \label{tab:modules}
  \small
  \begin{tabularx}{\textwidth}{@{}>{\raggedright\arraybackslash}p{2.7cm} R R@{}}
    \toprule
    \sffamily\bfseries\color{navy}Module &
    \sffamily\bfseries\color{navy}Rôle &
    \sffamily\bfseries\color{navy}Entrée $\rightarrow$ sortie \\
    \midrule
    \module{parser.py}    & Extraction du PDF, découpage par article, métadonnées de hiérarchie
                          & PDF source $\rightarrow$ corpus JSONL \\
    \module{database.py}  & Calcul des embeddings, construction de l'index dense
                          & JSONL $\rightarrow$ collection Chroma \\
    \module{retriever.py} & Recherche dense et lexicale, fusion RRF, renvoi par numéro explicite
                          & Question $\rightarrow$ articles + score \\
    \module{generator.py} & Prompt ancré, appel au modèle local
                          & Question + articles $\rightarrow$ réponse \\
    \module{app.py}       & Orchestration, abstention, contrôle des citations
                          & Question $\rightarrow$ réponse affichée \\
    \bottomrule
  \end{tabularx}
\end{table}

Chaque composant est isolé derrière une fonction simple, ce qui permet de remplacer une brique
sans toucher aux autres : changer de modèle de langage n'affecte que \module{generator.py} ;
passer du dense seul à l'hybride n'affecte que \module{retriever.py}. Cette contrainte est
délibérée, et sa justification est l'évaluation : elle permet de ne faire varier qu'une chose à la
fois. Seul le modèle d'embeddings fait exception --- il est partagé entre \module{database.py} et
\module{retriever.py}, pour la raison exposée plus haut.

\section{Choix technologiques}

\subsection{Deux filtres derrière chaque décision}

La quasi-totalité des choix découle de deux contraintes, appliquées dans cet ordre.

\begin{enumerate}
  \item \textbf{Souveraineté $\rightarrow$ tout doit s'exécuter en local.} Ce seul critère élimine
        les embeddings par interface distante, les services de reclassement commerciaux et les
        modèles de génération propriétaires.
  \item \textbf{Corpus borné et durée courte $\rightarrow$ optimiser la maîtrise, pas le passage à
        l'échelle.} À l'échelle de quelques centaines d'articles, toutes les bases vectorielles
        répondent instantanément : la performance brute n'est jamais le facteur décisif, et
        l'invoquer pour trancher aurait été un faux argument.
\end{enumerate}

\subsection{Décisions par composant}

Le tableau~\ref{tab:techno} récapitule les choix retenus couche par couche. Dans la majorité des
cas, la justification se ramène à l'un des deux filtres ci-dessus.

\begin{longtable}{@{}p{2.5cm} p{2.7cm} p{\dimexpr\textwidth-6.4cm\relax}@{}}
\caption{Choix technologiques et justifications}\label{tab:techno}\\
\toprule
\sffamily\bfseries\color{navy}Couche &
\sffamily\bfseries\color{navy}Choix &
\sffamily\bfseries\color{navy}Justification \\
\midrule
\endfirsthead
\toprule
\sffamily\bfseries\color{navy}Couche &
\sffamily\bfseries\color{navy}Choix &
\sffamily\bfseries\color{navy}Justification \\
\midrule
\endhead
\bottomrule
\endfoot
Extraction PDF & PyMuPDF & Rapide, fidélité Unicode fiable. Surtout : contrôle direct sur la façon dont la loi est découpée, là où un parseur automatique masque ce contrôle --- or c'est précisément ce découpage qui devait être vérifiable. \\[4pt]
Embeddings & BGE-M3 & La souveraineté exclut les embeddings distants. Parmi les modèles locaux : bonne couverture français et arabe, contexte long (8192 jetons), donc aucun article tronqué à l'indexation. \\[4pt]
Base vectorielle & Chroma & À cette échelle, la performance est hors sujet ; l'expérience de développement a donc primé. Chroma est embarquée (aucun serveur à administrer), persiste sur disque et filtre par métadonnées. \\[4pt]
Recherche lexicale & BM25 (\module{rank\_bm25}) & Les requêtes juridiques sont riches en jetons exacts --- numéros d'articles, «~Loi~65-99~» --- précisément là où le lexical bat les embeddings. \\[4pt]
Exécution du LLM & Ollama & Une commande télécharge et sert un modèle quantifié derrière une interface locale propre, sans qu'aucune requête ne sorte de la machine. \\[4pt]
Modèle de génération & Mistral~7B & La souveraineté exclut les modèles frontières distants. La classe 7B est le point d'équilibre : suffisante en suivi d'instructions, exécutable sur du matériel étudiant. Mistral pour la qualité de son français. \\[4pt]
Cadriciel RAG & \textbf{aucun} & Pour un corpus borné, la boucle RAG tient en environ 150 lignes. L'écrire soi-même permet de comprendre chaque étape, de déboguer directement, et d'expliquer son propre pipeline en soutenance. \\[4pt]
Format du corpus & JSONL & Un article par ligne : lisible, métadonnées imbriquées, comparable ligne à ligne sous Git, lecture en flux. \\
\end{longtable}

\subsection{Les arbitrages réellement discutables}

Il serait malhonnête de présenter tous ces choix comme évidents. Trois d'entre eux restent
ouverts, et sont énoncés ici pour ce qu'ils sont.

\begin{itemize}
  \item \textbf{BGE-M3 contre multilingual-e5.} Appel serré : e5 est plus léger et suffirait
        vraisemblablement. BGE-M3 l'emporte sur la largeur multilingue et la longueur de contexte,
        deux propriétés utiles pour l'extension envisagée --- donc un pari sur la suite, pas une
        supériorité mesurée.
  \item \textbf{Sans cadriciel contre LlamaIndex.} L'arbitrage le plus discutable. Un cadriciel se
        justifierait si le projet grossissait : sources multiples, agents, routage. Le refuser est
        un choix pédagogique assumé, dont le coût serait payé par le prochain développeur.
  \item \textbf{Dense seul contre hybride.} L'hybride est peu coûteux et le texte juridique
        favorise la correspondance exacte ; l'argument théorique lui est favorable. La position
        honnête est cependant \emph{«~je l'ai mesuré, voici l'écart~»} --- et le
        chapitre~\ref{chap:evaluation} montre que cet écart dépend surtout des questions qu'on
        pose. Mesuré deux fois sur deux jeux de questions différents, il change de signe.
\end{itemize}

\section*{Conclusion}
\addcontentsline{toc}{section}{Conclusion}

L'architecture retenue place délibérément le garde-fou d'abstention hors du modèle et maintient
la recherche séparable de la génération. Ces deux décisions, prises tôt, ont conditionné tout ce
qui a pu être débogué et mesuré ensuite : sans la première, l'abstention aurait dépendu du modèle
lui-même ; sans la seconde, aucune des deux causes possibles d'une mauvaise réponse n'aurait pu
être distinguée de l'autre. Le chapitre suivant décrit leur mise en œuvre concrète.
